\documentclass[10pt,a4paper]{amsart}
\usepackage[margin=19mm]{geometry}
\usepackage{amsmath,amssymb,mathtools}
\usepackage[scr=rsfs,cal=euler]{mathalfa}
\usepackage{xfrac}
\usepackage[unicode,hidelinks,bookmarks=false]{hyperref}
\newcommand{\ie}{i.e.\ }
\newcommand{\cf}{cf.\ }
\newcommand{\resp}{respectively\ }
\newcommand{\Subsection}{\subsection}
\providecommand{\nobreakdash}{\nobreak\hbox{-}}
\setlength{\emergencystretch}{2em}
\multlinegap0pt
\usepackage{amssymb}
\usepackage{enumerate}
\usepackage{tikz}
\usepackage{mathtools}
\usepackage[shortlabels]{enumitem}
\newtheorem{prop}{Proposition}
\newcommand{\lie}[1]{\mathfrak{#1}}
\title{Graded embeddings, root generated subalgebras and $\pi$-systems for quasisimple Kac-Moody superalgebras}
\author{Irfan Habib, Deniz Kus, Chaithra Pilakkat}
\date{Source: arXiv:2512.19222v1 (CC BY 4.0)}
\begin{document}
\maketitle

\noindent\emph{Source and licence.} Graded embeddings, root generated subalgebras and $\pi$-systems for quasisimple Kac-Moody superalgebras. Version arXiv:2512.19222v1 (CC BY 4.0), distributed by arXiv under the Creative Commons Attribution 4.0 International licence, http://creativecommons.org/licenses/by/4.0/, as recorded in the arXiv licence metadata for this exact version and shown on the abstract page https://arxiv.org/abs/2512.19222v1. Full licence notice: \texttt{licenses/arxiv-2512.19222-CC-BY-4.0.txt}.\par\smallskip

\noindent\textit{Editorial scope.} Counted unit: the article's prop unbrokenrs (source0.tex lines 483--490). The article's complete original proof of it is delivered with it (source0.tex lines 491--506, one \texttt{proof} environment of 16 lines). No mathematics is written, repaired or re-derived: the statement and the proof are the article's own lines, in the article's own notation, and no retained line is substituted or re-flowed.  No further source-local context is needed: every cross-reference of every retained line resolves inside the two delivered spans.  Nothing was omitted from any retained span, so the package carries no omission record.  No retained line contains a citation, so the wrapper carries no bibliography and no bibliography entry.  No retained line contains a figure environment, an image-inclusion command or drawing code, so no image is delivered and the package declares no asset dependency.  Editorial formatting only, in addition: an AMS \texttt{amsart} preamble that restates the article's own declarations and macros used by the retained lines, namely the environments \texttt{prop} on separate counters, together with the standard packages those lines need.  The retained environments print in this document as Proposition 1 in order of appearance, whereas the article prints the same items with the numbering of its own full document.  The complete native source of the article as distributed in the arXiv e-print for this exact version is bundled unchanged as \texttt{sources/191454/source0.tex}.
\begin{prop}\label{unbrokenrs}
Let $\alpha,\beta\in \Delta$ be two roots of a Kac-Moody superalgebra such that $\alpha$ is real.

\begin{enumerate}
    \item  If $\alpha$ is non-isotropic, then there exist non-negative integers $p$ and $q$ related by the equation $p-q=\beta(h_\alpha)$ such that $\beta+k\alpha\in \Delta\cup \{0\}$ if and only if $-p\le k\le q.$ In particular, $s_\alpha$ reverses the root string $\mathbb{S}(\beta,\alpha)$. \vspace{0,1cm}
    \item If $\alpha$ is isotropic, then $|\mathbb{S}(\beta,\alpha)|$ is finite.
\end{enumerate}
\end{prop}

\begin{proof}
 We start proving part (1). The result is easily checked if $\beta$ and $\alpha$ are scalar multiple of each other. Otherwise, consider the subspace $U=\sum_{k\in \mathbb{Z}} \lie g_{\beta+k\alpha}$ of $\lie g.$ Since $\alpha$ is non-isotropic, the subalgebra $\lie g(\alpha)$ 
    is isomorphic to either $\lie{sl}_2$ or $\lie{osp}(1,2).$ Let $v$ be a root vector in $U.$ Since $\mathrm{ad}_{x_\alpha^\pm}$ acts locally nilpotently, the $\lie g{(\alpha)}$ submodule of $U$ generated by $v$ is finite-dimensional. Thus, by the representation theory of $\lie{sl}_2$ or $\lie{osp}(1,2)$, we have that $U$ is a sum of finite-dimensional irreducible $\lie g{(\alpha)}$-modules. 
    Since the eigenvalues of $h_\alpha$ in $U$ are of the form $\beta(h_\alpha)+2k$, it follows that every irreducible module that occur in $U$ contains a non-zero vector in $\lie g_{\beta+t\alpha}$ where $$t=\begin{cases}
        \frac{-\beta(h_\alpha)}{2}& \text{ if } \beta(h_\alpha) \text{ is even},\\
        \frac{1-\beta(h_\alpha)}{2}& \text{ if } \beta(h_\alpha) \text{ is odd}.
    \end{cases}$$
    Therefore, $U$ is equal to the $\lie g{(\alpha)}$ submodule generated by the root space $\lie g_{\beta+t\alpha}$. In particular, $U$ is finite-dimensional and the claim follows from the representation theory of $\lie g(\alpha).$
    
  Now we prove part (2). If $\alpha$ and $\beta$ are scalar multiples of each other, the statement is clear. Since $\alpha$ is real and odd we must have $\alpha\in\Sigma$ for a base $\Sigma$. Moreover, 
\begin{equation}\label{dreir}[x^+_{\alpha},[x_{\alpha}^+,y]]=0,\ \ y\in \lie g.\end{equation} We first assume that $\beta\in C^+_\Sigma$, i.e., $\beta=\sum_{\gamma\in \Sigma}a_{\gamma}\gamma$ where $a_{\gamma}\geq 0$. If $k\in \mathbb{Z}_+$ is such that $\beta+k\alpha\in \Delta$, we can express any element in $\mathfrak{g}_{\beta+k\alpha}$ as a linear combination of Lie words in the root vectors corresponding to the base $\Sigma$. However, by \eqref{dreir} we must have 
$$\mathfrak{g}_{\beta+k\alpha}\neq 0 \implies k\leq \sum_{\gamma\in \Sigma\setminus\{\alpha\}}a_\gamma-a_\alpha.$$
Moreover, if $t > a_\alpha$, we claim that $\beta - t\alpha$ is never a root. Indeed, we have
$$\beta - t\alpha = \sum_{\gamma \in \Sigma \setminus \{\alpha\}} a_{\gamma} \gamma + (a_\alpha - t)\alpha.$$
Since $\beta$ is not a scalar multiple of $\alpha$, there exists some $\gamma \in \Sigma \setminus \{\alpha\}$ with $a_{\gamma} > 0$. Since $t > a_\alpha$, we must have $\beta - t\alpha\notin \Delta$. The proof  for the case $\beta\in -C^+_\Sigma$ is similar.
\end{proof}

\end{document}
